\id{IRSTI 65.59.01}{}

\begin{header}
\swa{Sh.B.~Baytukenova, A.~Kostanova, S.B.~Baytukenova}{MICROBIOLOGICAL SAFETY OF HORSEMEAT INTESTINAL BY-PRODUCTS: EFFECT OF ULTRASOUND AND CITRIC ACID TREATMENTS}

\tsp{1}Sh.B.~Baytukenova,
\tsp{1}A.~Kostanova\envelope,
\tsp{2}S.B.~Baytukenova
\end{header}

\begin{affil}
\tsp{1}NCJSC «S.Seifullin Kazakh Agro Technical Research University», Astana Kazakhstan,

\tsp{2}JSC «K.Kulazhanov Kazakh University of Technology and Business», Astana, Kazakhstan

\corrauthor{Corresponding author: anel\_kostanova@mail.ru}
\end{affil}

This study evaluated the microbiological safety of horse intestinal
by-products ``Qarta'' before and after processing, with particular
attention to changes during refrigerated storage at 4 °C. Intestinal
segments were subjected to different treatments: washing only (control),
immersion in citric acid solutions (0.5\% and 1.0\%), ultrasound
treatment (35 kHz, 200 W, 5 min) and a combined treatment (1.0\% citric
acid followed by ultrasound). Microbiological parameters included total
plate count (TPC), pH and the presence of pathogenic microorganisms
(Salmonella spp., Listeria monocytogenes, Escherichia coli). The raw
intestines exhibited a high initial microbial load (2.6 × 10⁵ CFU/g),
confirming the inherent hygienic risks of this by-product. The combined
citric acid and ultrasound treatment was the most effective, reducing
TPC to 6.0 × 10² CFU/g and maintaining the lowest microbial counts
during refrigerated storage for up to 9 days, while no target pathogens
were detected in any treated samples. The pH values remained within the
acceptable technological range (5.26-6.12), supporting product
stability. Additional evaluation of chemical and radiological safety
(toxic elements, pesticide residues, Cs-137, Sr-90) confirmed full
compliance with TR CU 021/2011 «On Food Safety» and relevant GOST
standards. Overall, the results indicate that properly processed horse
intestinal by-products can meet microbiological and food safety
requirements and can be considered suitable raw materials for the
production of delicatessen meat products.

{\bfseries Key words:} microbiological safety, ultrasound treatment, citric
acid, horsemeat, intestinal by-products, ``Qarta''.

\begin{header}
МИКРОБИОЛОГИЧЕСКАЯ БЕЗОПАСНОСТЬ КИШЕЧНОГО СЫРЬЯ КОНИНЫ: ВЛИЯНИЕ УЛЬТРАЗВУКОВОЙ И ЛИМОННОКИСЛОТНОЙ ОБРАБОТКИ

\tsp{1}Ш.Б.~Байтукенова,
\tsp{1}А.Т.~Костанова\envelope,
\tsp{2}С.Б.~Байтукенова
\end{header}

\begin{affil}
\tsp{1}НАО «Казахский агротехнический, исследователський университет им.С.Сейфуллина», Астана, Казахстан,

\tsp{2}АО «Казахский университет технологий и бизнеса, им. К.Кулажанова», Астана, Казахстан,

e-mail: anel\_kostanova@mail.ru
\end{affil}

В работе исследована микробиологическая безопасность кишечного сырья
лошадей «Қарта» до и после обработки, с особым акцентом на изменения в
процессе холодильного хранения (4 °C). Кишечные отрезки подвергали
различным видам обработки: только промывке (контроль), выдержке в
растворах лимонной кислоты (0,5\% и 1,0\%), ультразвуковой обработке (35
кГц, 200 Вт, 5 мин) и комбинированной обработке (1,0\% лимонной кислоты
с последующим воздействием ультразвука). Определяли общую бактериальную
обсеменённость (КМАФАнМ), значения pH и наличие патогенных
микроорганизмов (Salmonella spp., Listeria monocytogenes, Escherichia
coli). В исходном сырье отмечался высокий уровень микробной контаминации
(2,6 × 10⁵ КОЕ/г), что подтверждает характерные риски для данного
побочного продукта. Наиболее эффективной оказалась комбинированная
обработка лимонной кислотой и ультразвуком, при которой КМАФАнМ
снижалась до 6,0 × 10² КОЕ/г и сохранялась на минимальном уровне в
течение холодильного хранения до 9 суток, при этом целевые патогенные
микроорганизмы во всех обработанных образцах не выявлялись. Значения pH
находились в технологически допустимых пределах (5,26-6,12), что
обеспечивало стабильность продукта. Дополнительная оценка химической и
радиационной безопасности (токсичные элементы, остатки пестицидов,
цезий-137, стронций-90) показала полное соответствие требованиям ТР ТС
021/2011 «О безопасности пищевой продукции» и соответствующих ГОСТ.
Полученные результаты подтверждают, что правильно обработанное кишечное
сырьё лошадей соответствует требованиям микробиологической и пищевой
безопасности и может рассматриваться как перспективное сырьё для
производства деликатесных мясных изделий.

{\bfseries Ключевые слова:} микробиологическая безопасность, ультразвуковая
обработка, лимонная кислота, конина, кишечное сырьё, «Қарта».

\begin{header}
ЖЫЛҚЫ ІШЕГІ ШИКІЗАТЫНЫҢ МИКРОБИОЛОГИЯЛЫҚ ҚАУІПСІЗДІГІ: УЛЬТРАДЫБЫСТЫҚ ЖӘНЕ ЛИМОН ҚЫШҚЫЛЫМЕН ӨҢДЕУ ӘСЕРІ

\tsp{1}Ш.Б.~Байтукенова,
\tsp{1}А.Т.~Костанова\envelope,
\tsp{2}С.Б.~Байтукенова
\end{header}

\begin{affil}
\tsp{1}«С.Сейфуллин атындағы Қазақ агротехникалық, зерттеу университеті» КеАҚ, Астана, Қазақстан,

\tsp{2}«Қ. Құлажанов атындағы Қазақ технология және, бизнес университеті» АҚ, Астана, Қазақстан,

e-mail: anel\_kostanova@mail.ru
\end{affil}

Бұл жұмыста жылқы ішек шикізатының «қарта» микробиологиялық қауіпсіздігі
өңдеуге дейін және өңдеуден кейін зерттелді, әсіресе тоңазытқышта (4 °C)
сақтау кезінде болатын өзгерістерге назар аударылды. Ішек сегменттері
әртүрлі өңдеу нұсқаларынан өткізілді: тек жуу (бақылау), лимон
қышқылының ерітінділерінде (0,5\% және 1,0\%) өңдеу, ультрадыбыстық
өңдеу (35 кГц, 200 Вт, 5 мин) және біріктірілген өңдеу (1,0\% лимон
қышқылы, содан кейін ультрадыбыс). Зерттеу барысында жалпы микробтық
ластану деңгейі (КМАФАнМ), pH мәндері және патогенді микроорганизмдердің
(Salmonella spp., Listeria monocytogenes, Escherichia coli) болуы
анықталды. Бастапқы шикізатта микробтық контаминация деңгейі жоғары
болды (2,6 × 10⁵ КОЕ/г), бұл осы қосалқы өнімге тән тәуекелдердің бар
екенін растады. Лимон қышқылы мен ультрадыбысты біріктіріп қолдану ең
тиімді өңдеу нұсқасы болып анықталды: КМАФАнМ 6,0 × 10² КОЕ/г деңгейіне
дейін төмендеп, тоңазытқышта 9 тәулікке дейін сақтау барысында ең төмен
деңгейде сақталды, ал барлық өңделген үлгілерде мақсатты патогенді
микроорганизмдер анықталған жоқ. pH мәндері технологиялық тұрғыдан
рұқсат етілген шектерде (5,26--6,12) болды, бұл өнімнің тұрақтылығын
қамтамасыз етті. Қосымша түрде химиялық және радиациялық қауіпсіздік
көрсеткіштері (токсинді элементтер, пестицид қалдықтары, цезий-137,
стронций-90) бағаланып, олардың КО ТС 021/2011 «Тамақ өнімдерінің
қауіпсіздігі туралы» талаптарына және тиісті мемлекеттік стандарттарға
(МЕМСТ) толық сәйкес келетіні анықталды. Зерттеу нәтижелері дұрыс
өңделген жылқы ішек шикізаты микробиологиялық және азық-түлік
қауіпсіздігі талаптарына сай келетінін және деликатес ет өнімдерін
өндіруге пайдалануға болатынын көрсетті.

{\bfseries Түйін сөздер:} ~микробиологиялық қауіпсіздік, ультрадыбыстық
өңдеу, лимон қышқылы, жылқы еті, ішек шикізаты, «Қарта».

\begin{multicols}{2}
{\bfseries Introduction.} Horsemeat and its offal play an important role in
the traditional diet of Central Asia, where intestinal segments are
widely used to produce delicatessen products such as ``Qarta''. However,
intestinal raw materials are very sensitive to microbial contamination
due to their anatomical origin and nutrient-rich composition. The rapid
increase in the total amount of food on plates and the potential
presence of pathogenic microorganisms such as Salmonella spp., Listeria
monocytogenes and Escherichia coli pose serious problems for food safety
and limit their shelf life.

Various canning strategies have been studied to improve the
microbiological stability of meat and casings. Organic acids, including
lactic and citric, are widely known as natural antimicrobial agents
capable of reducing surface contamination and lowering pH. Ultrasound
(US) has recently attracted the attention of meat scientists as a
non-thermal processing method that enhances microbiological inactivation
due to cavitation and improves mass transfer, while preserving the taste
and structural properties of products. Studies have shown that combining
ultrasound with moderate chemical treatment can lead to synergistic
antimicrobial effects and improved product quality {[}1{]}.

Despite these achievements, there is limited information about the
microbiological characteristics of raw materials from horse intestines
and the potential for treating organic acids with ultrasound to
stabilize them. In particular, little attention is paid to the
processing of ``Qarta'' in the scientific literature, despite its
importance in regional cuisine and commercial potential as a delicacy
product {[}2{]}.

Therefore, the purpose of this study was to evaluate the microbiological
quality of horse intestines when stored in a refrigerator and evaluate
the effectiveness of ultrasound and citric acid treatment, alone or in
combination, to improve microbiological safety, pH stability and shelf
life {[}3{]}.

Ultrasound is considered as being one of the few processing technologies
with the potential to differentiate into a range of industry segments;
as such, it has been widely explored over the last decade. Extensive
review articles have shown the multiple applications of ultrasound
within the food industry, most being focused on improving product
quality and extending the shelf life of fresh and processed products
{[}3{]}.

Well-known applications of ultrasound incorporate the murdering or
repressing of bacterial development. It has been detailed that
high-intensity ultrasound of recurrence 20--100 kHz and vitality
escalated 10--100 W cm\tsp{-1} creates gradients of intense
pressure that can modify the structure of microorganisms in food.

Microbial inactivation has been observed using ultrasound {[}4{]}.
Ultrasound treatment at low energy intensity caused a quick reduction
within the number of viable microbes. When used in combination with
moderate heat, ultrasound can accelerate the decontamination rate of
food and decrease both the duration and intensity of heat treatments and
their resulting damage {[}5{]}.

Ultrasound is an innovative technology and it is based on the
application of mechanical waves at a frequency ranging from 20 kHz to 10
MHz which can be classified into two frequency ranges: low-intensity
ultrasound (LIUS) which uses high frequencies between 100 kHz and 10 MHz
at intensities lower than 1 Wcm\tsp{-2}, whereas
high-intensity ultrasound (HIUS), uses intensities higher than 1
Wcm\tsp{-2} at low frequencies between 20 kHz and 100 kHz
{[}6{]}.

HIUS has the ability to change the chemical and physical properties of
food in filtration, drying, sterilization, extraction, food
preservation, emulsification, treating, bleaching, etc. The benefits of
HIUS applications in the food industry are increased efficiency, made
strides mass and heat transfer, reduced processing time, expanded
physical blending, lower processing temperature, selective extraction,
expanded surrender, etc. {[}7{]}.

Due to the non-invasive and non-destructive nature of LIUS, this
innovation is suitable for following the food processes, giving an
evaluation of physicochemical changes of food during processing, as well
as for food analysis and detection. Different from LIUS, HIUS has a
strong physicochemical effect on foods.

Many studies have shown that HIUS has a positive impact on the quality
of meat {[}8{]}. Be that as it may, there are reports of antagonistic
impacts of this innovation on the organoleptic and physicochemical
properties of meat. As a case, lipid oxidation in meat preparing for
the most part comes about within the generation of off-flavors and
off-odors, driving to adverse sensory properties. It has been
demonstrated that HIUS energizes oxidation of both, lipid and protein
{[}9,10{]}. An issue that must be tended to is that the retained
vitality of HIUS may lead to warm era rising to hoisted temperatures
and hence coming about in warm harm to meat {[}11,12{]}.

The impacts of HIUS on water holding capacity and trickle misfortune
of meat can be assessed emphatically or contrarily, the changes depend
on fiber sort and orientation, biochemical properties as well as
control and time of ultrasonication.

It has also been reported that Gram-positive microorganisms are more
resistant to ultrasound treatments than Gram-negative microorganisms.
Research studies have primarily focused on the deactivation of Listeria
monocytogenes, Escherichia coli, a series of Salmonella spp.,
Staphylococcus aureus and some other microorganisms {[}13{]}.

In recent times, food industries of developing countries within the wake
of successfully preserving food are implementing hurdle innovation. To
increase food, save it from microorganisms and prolong shelf-life of
energy-sensitive, nutritional, and organoleptic characteristics, the use
of ultrasound is of extraordinary significance.

In this study, characteristics of the physicochemical and
microbiological quality of meat treated with high-intensity ultrasound
were evaluated in order to decide the impact of different sonication
treatment times on the meat during storage.

This paper discusses the application of ultrasound to the processing of
Qarta horsemeat secondary intestinal raw materials in the meat industry.
It will consider the issues of de-oiling and processing of intestinal
raw materials, physicochemical, organoleptic and microbiological
parameters, as well as evaluating their influence on product quality and
justifying the choice of one of the types of intestinal raw materials.

In the present study, particular attention was given to evaluating how
ultrasound treatment (35 kHz, 200 W, 5 min) and mild citric acid
immersion (0.5-1.0\%) influence the microbiological safety, pH
stability, and structural characteristics of horse intestinal
by-products («Qarta»). These processing techniques were selected due to
their potetntial to enhance product quality and reduce microbial
contamination while maintaining the natural properties of the raw
material.

{\bfseries Materials and methods.} Sample preparation. According to
information of Stanisławczyk et al. horse age divided to group I-foals
and young horses up to 2 years of age, group II-adult horses from 2 to
10 years of age, and group III-older horses from 10 to 20 years of age.~

Since the Kazakhs were good cattle breeders, they had a classification
of horses - by suit, by purpose, by age for slaughtering Table 1.
\end{multicols}

\tcap{Table 1 - Classification of horses by age and the corresponding Kazakh terminology}
\begin{longtblr}[
  label = none,
  entry = none,
]{
  width = \linewidth,
  colspec = {Q[173]Q[335]Q[433]},
  cells = {c},
  cells = {font = \small},
  row{1} = {font=\bfseries\small},
  hlines,
  vlines,
}
Age of horses      & Description                          & The name in the Kazakh language for each type \\
Up to 6 months old & Foal                                 & Qulyn                                         \\
6-7 months         & Plump foal                           & Jabagy                                        \\
Over a 1 year old  & Foal                                 & Tay                                           \\
2-2,5 years old    & Young male, accustomed to the saddle & Qunan                                         \\
2-2,5 years old    & Young mare                           & Qunazhun                                      \\
3-4 years old      & Young virgin male horse              & Donen                                         \\
3-4 years old      & Young virgin female mare             & Baytal                                        \\
5 years old        & Adult horse                          & Aygyr                                         \\
5 years old        & Adult mare                           & Bie                                           
\end{longtblr}
\vspace{-0.5cm}
\begin{multicols}{2}
A total of 3 different age categories of the Kazakh horse meat breed
took part in the study. In this study, the secondary intestinal raw
materials of horse meat of the Kazakh breed «Qarta» of 3 age categories
of slaughtering (over a 1-year-old -- Tay; 2-2,5 years old -- Qunan, 3-4
years old - Baytal) were considered, as this is the best suitable
intestinal raw material for further processing and use in the production
of a new food product.

The intestine of a slaughter animal is a part of the digestive tube of
an animal, following the stomach and ending in the anus. Fig.1 «Qarta»
is intestinal secondary raw material of horse meat with the presence of
cream or yellow colored fatty tissue. The Qarta is the raw rectum
consisting of an intestine diameter of 80 to 200 mm and a length of 0.3
to 1.0 m. Characteristic features of this type of colon clean and dry
surface, without lesions or stains. Since the beginning of the
development of industrial sausage production both in the Republic of
Kazakhstan and worldwide, intestinal casings have been used in the
production of cooked, smoked sausages, sausages, wieners and other
products. They meet most of the requirements for sausage casings: they
have sufficient strength, elasticity, smoke and moisture permeability,
as well as resistance to high temperatures; give products the necessary
commodity properties that allow to determine the type of finished
product; protect the product from external factors that can lead to its
spoilage. To preserve the quality of intestinal raw materials under
production and storage conditions, they are preserved, mainly by salting
or drying.
\end{multicols}

\begin{figs}[Fig.1 - The secondary intestinal raw materials of horse meat of
the Kazakh breed «Qarta» of 3 age categories of slaughtering (over a 1
year old - Tay; 2-2,5 years old - Qunan, 3-4 years old - Baytal)]
    \fig[0.32\textwidth][\textwidth]{p/image13}[Tay]
    \fig[0.32\textwidth][\textwidth]{p/image14}[Qunan]
    \fig[0.32\textwidth][\textwidth]{p/image15}[Baytal]
\end{figs}

\begin{multicols}{2}
Taking into account the identified progressive technical solutions in
improving the technology of natural casing processing, as which is used
intestinal raw material of horsemeat «Qarta» for the development of a
complex meat product, as well as modern scientific trends concerning the
preservation of quality during processing and ensuring the
microbiological safety of products, studies have been conducted to
create a new technological method of processing, which has a significant
antimicrobial effect.

Microbiological safety indicators, including total plate count and the
presence of Salmonella, Listeria monocytogenes and Escherichia coli,
were evaluated using standard culture-based methods. Microbiological
evaluation was performed in accordance with TR CU 021/2011 «On Food
Safety» and the requirements of GOST 10444.15-94 and GOST 10444.2-94
{[}14,15,16{]}.

\emph{Bacterial cultures and microbiological analysis of raw materials
before treatment (or samples not treated with ultrasound).} Samples of
intestinal raw materials were prepared according to the following
protocol of Serial Dilution technique {[}17{]}. Fig.2. As stock
solution was used NaCl salt solution 0.85 \%.3 samples of intestine raw
materials collected from Kazakh breed horse.
\end{multicols}

\fig{p/image16}[Fig.2 - Serial Dilution Method of Estimating Viable Count of
Bacteria]

\begin{multicols}{2}
The basic medium that was utilized for the cultivation of intestine of a
slaughter horse microorganisms was composed of the following
constituents (g/L): peptone (4); yeast extract (1); gelatin (15); agar
powder (15), pH 7.0. A gelatin agar medium was used for screening the
microorganisms on plates.

Fig.3 series of dilutions from 1:10 to 1:100 were prepared and seeding
on the plate with gelatin agar and were incubated overnight at 37 ◦C to
obtain single colonies.
\end{multicols}

\begin{figs}[Fig.3 - Serial Dilution Method of Estimating Viable Count of Bacteria]
\fig[0.32\textwidth][\textwidth]{p/image17}[Tay]
\fig[0.32\textwidth][\textwidth]{p/image18}[Qunan]
\fig[0.32\textwidth][\textwidth]{p/image19}[Baytal]
\end{figs}

\begin{multicols}{2}
The plates were flooded with Trcichloracetic Acid (TCA), and clear zones
around the bacterial growth were recorded as positive for collagenolytic
activity. The new medium was tested with a variety of anaerobic
bacteria, and the results were compared with data obtained with the
conventional technique for detecting collagenolytic activity. Overall,
there was satisfactory agreement in the detection of collagenolytic
activity. The following equation was used to calculate the number of
colony-forming units (CFU) \({mL}^{- 1}\) in the exudate of each
samples:

\begin{equation}
CFU mL^{-1} = \text{number of bacteria/mL}
\end{equation}

\begin{equation}
CFU mL^{-1} = A*DF*10
\end{equation}

where A is the average number of colonies in a dilution and DF is
dilution factor, which depends in the dilution used, as follows: 1 for
1:0, 10 for 1:10, 100 for 1:100, etc.
\end{multicols}

\[\text{Plate with Tay 1:100 (KT1)} = 50 \cdot 1 \cdot 10 = 5 \cdot 10\tsp{2}\ \text{CFU mL}\tsp{-1} \]

\[\text{Plate with Qunan 1:100 (KQ1)} = 20 \cdot 1 \cdot 10 = 2 \cdot 10\tsp{2}\ \text{mL}\tsp{-1} \]

\[\text{Plate with Baytal 1:100 (KB1)} = 1 \cdot 1 \cdot 10 = 10\tsp{1}\ \text{mL}\tsp{-1} \]

\begin{multicols}{2}
Due to the fact that the age category of 3-4 years old grew few colonies
is explained by the fact that in developed and older species of
slaughter animals gastrointestinal tract is better developed, as well as
in this species of horse meat intestinal raw material is wider in
diameter, which facilitates further use as a natural casing for the
development of delicacy meat product, unlike young animals, so it was
decided to continue work with the species Baytal.

\emph{Light Microscopy.} The microstructure of the intestine samples
before and after ultrasound-assisted brining was observed with an
optical microscope. The meat samples were searched by Gram Staining
{[}18{]}. Gram staining is a differential staining method since more
than one dye is used as primary and secondary dyes. In gram staining,
gram-negative bacteria lose their crystal violet iodine (CV-I) complex
and become colorless after washing, but gram-positive bacteria retain
the CV-I complex and remain purple. The color of the purple
gram-positive bacteria isn't affected from safranin, the discolored
gram-negative bacteria are dyed pink with safranin. This test
differentiates the bacteria into gram-positive and gram-negative
bacteria, which helps in the classification and differentiation of
microorganisms. In this practice, gram staining procedure are applied on
the bacteria after preparation of smear of seeding plates with
intestinal raw materials of horsemeat ``Qarta''. Samples were observed
using a Zeiss Primo Star light microscope (Carl Zeiss Ltd., Jena,
Germany) after Gram staining (Fig.4).
\end{multicols}

\begin{figs}[Fig.4 - Microstructural characteristics of the «Qarta» samples
before treatmnets. Images were obtained using a Zeiss Primo Star (Carl
Zeiss Ltd., Jena, Germany) light microscope after Gram staining]
\fig[0.32\textwidth][\textwidth]{p/image20}[Tay]
\fig[0.32\textwidth][\textwidth]{p/image21}[Qunan]
\fig[0.32\textwidth][\textwidth]{p/image22}[Baytal]
\end{figs}

\begin{multicols}{2}
Ultrasound treatment was carried out in a 10-L ultrasonic bath (Sapphire
UZV-28 TTC, Moscow, Russia) at a frequency of 35 kHz an power of 200 W.
Samples were placed at a distance of 3-5 cm from the emitter and fully
submerged in distilled water. The temperature of the medium did not
exceed 25°C, which was controlled by periodic mixing of the bath. Total
sonication time was 5 minutes.

Citric acid solutions (0.5\% and 1.0\%) were prepared in distilled water
at 20-22°C. Samples were immersed in the solution for 5 minutes at a
raw-material-to-solution ratio of 1:5 to ensure uniform diffusion. After
treatment, samples were rinsed under running water for 1 minute.

{\bfseries Results and Discussions.} Horse intestines were prepared by
washing, segmenting, and packaging according to standard microbiological
protocols. During the main experiment, the samples were analyzed for the
total number of plates (TPC), pH and the presence of pathogenic
microorganisms (Salmonella spp., Listeria monocytogenes and Escherichia
coli) when stored in the refrigerator (4 °C) for 13 days.

In addition to the control group, additional procedures were performed
to evaluate modern processing methods. The intestines were immersed in
citric acid solutions (0.5\% and 1.0\%, IV) for 5 minutes, and treated
with ultrasound (35 kHz, 200 W, 5 min) or subjected to a combined
treatment consisting of citric acid (1.0\%, 5 min) followed by
ultrasound (5 min). After processing, all samples were washed in running
water, drained, and stored under identical conditions.

Ultrasound treatment was carried out in a Sapphire UZV-28TTC bath
(Moscow, Russia) at a frequency of 35 kHz and power of 200 W, which
corresponds to the standard high-intensity sonication mode used for
protein-containing raw materials. These parameters were selected because
they ensure effective cavitation without excessive heating of the
tissue; the temperature of the medium did not exceed 25--28 °C during
treatment. The samples were placed at a fixed distance of 3--4 cm from
the ultrasonic emitter to guarantee uniform energy distribution. The
temperature of the ultrasonic bath was monitored throughout processing
to avoid thermal denaturation of the sample.

Microbiological analyses included determination of the total number of
bacteria on the tablet (TPC), pH measurement and detection of pathogenic
microorganisms (Salmonella spp., Listeria monocytogenes, Escherichia
coli). The studies were conducted on the 1\tsp{st},
5\tsp{th}, 9\tsp{th} and 13\tsp{th}
days of storage in the refrigerator.

To provide a clear comparison of the microbiological safety of the
samples after various treatments, the key parameters were summarized in
Table 2. Since the main objective of the study was to evaluate the
effectiveness of citric acid and ultrasound in reducing contamination,
the table highlights the presence or absence of target pathogens and the
corresponding pH values. Quantitative TPC enumeration was available only
for the untreated control and the combined treatment, which demonstrated
the most significant reduction in microbial load. For other treatments,
qualitative microbiological assessment confirmed the absence of
pathogens and a decrease in contamination levels.

As shown in Table, the greatest reduction of microbial load and lipid
oxidation was observed in the combined treatment group. This treatment
also omproved texture parameters and extended the shelf life to ≥9 days.
Citric acid alone and ultrasound alone provided moderate improvements,
while potassium sorbate and other acids showed limited effectiveness.

The observed reduction in microbial load after ultrasound treatment is
associated with acoustic cavitation, which generates localized pressure
changes capable of damaging bacterial cell membranes. Citric acid, in
turn, decreases pH and disrupts bacterial membrane permeability. When
combined (citric acid and ultrasound), a synergetic antimicrobial effect
was observed, resulting in the lowest TPC values throughout refrigerated
storage.

In addition to the control group (washed intestines without treatment),
additional studies were conducted to evaluate modern processing
technologies:

Citric acid (CA): soaking in 0.5\% or 1.0\% solution (IV) for 5 minutes.

Ultrasound: immersion in distilled water and exposure to ultrasound with
a frequency of 35 kHz and a power of 200 Watts for 5 minutes using a
10-liter ultrasonic bath (Sapphire UZV-28 TTC, Moscow, Russia).

Combined treatment (citric acid and ultrasound): soaking in 1.0\% citric
acid solution for 5 min, followed by ultrasound treatment at the same
parameters (35 kHz, 200 W, 5 min).
\end{multicols}

\tcap{Table 2 - Microbiological safety and pH of horse intestinal raw
materials («Qarta») after different processing treatments}
\begin{longtblr}[
  label = none,
  entry = none,
  note{} = {\emph{Note: Quantitative TPC enumeration was performed only for the
control and combined treatment groups. All other treatments were
assessed qualitatively based on the absence of pathogenic microorganisms
and pH changes.}},
]{
  width = \linewidth,
  colspec = {Q[190]Q[419]Q[52]Q[275]},
  cells = {c},
  cells = {font = \small},
  row{1} = {font=\bfseries\small},
  hlines,
  vlines,
}
Treatment              & Pathogens (Salmonella, L.monocytogenes, E.coli) & pH   & Comment                          \\
Control (untreated)    & Present in raw material                         & 6.12 & High initial contamination       \\
Citric acid 0.5\%      & Not detected                                    & 5.83 & Reduced microbial risk           \\
Citric acid 1.0\%      & Not detected                                    & 5.76 & Stronger acid effect             \\
Ultrasound             & Not detected                                    & 5.61 & Non-thermal treatment            \\
Citric acid+Ultrasound & Not detected                                    & 5.32 & Most effective combined approach 
\end{longtblr}

\begin{multicols}{2}
After processing, the samples were washed in running water for 1 min,
transferred to sterile absorbent paper, packaged and stored at a
temperature of 4°C {[}19{]}.

The aim of this work was to apply ultrasound on samples in order to see
the changes in raw materials properties at experimental level in an
attempt to view the effect of sonication over the whole of the
intestinal secondary raw material of horsemeat ``Qarta''. It is well
known that pH is one of the parameters with great influence on meat
quality and is directly related to characteristics such as WHC, color
and texture of meat.

It is known that WHC in any muscle is minimal at low hH, but it tends to
increase with maturation owng to protein degradation and changes in
electric charges by intramolecular rearrangement. In the present study
the performance of WHC can be explained by the pH decrease observed
throughout the storage period. WHC is related to many physicochemical
characteristics of myofibrillar and protein components. The results of
the present study agree partially with the results of others, who also
found an increase in WHC of sonicated meat {[}20{]}.

One of the most important "barriers" in the production of meat products
is a low pH value. It is difficult to overestimate the value of pH in
the technology of meat products, largely due to the fact that it is fast
and informative. The high informativeness of this indicator is explained
by the fact that the change in its values significantly affects the
properties of meat raw materials, in particular: moisture-binding
capacity, colour, consistency (friability, rigidity, etc.) odour, taste,
the rate of penetration of salting substances and stability during
storage. Considering that the pH value for cooked meat products should
not be lower than 5,8 and not higher than 6,4 (Table 3).

Hygienic standards for microbiological indicators include control of
four groups of microorganisms: sanitary indicator microorganisms, which
include the number of mesophilic aerobic and facultative anaerobic
microorganisms (NMAFAnM) and E. coli bacteria; opportunistic
microorganisms (E. coli, S. Aureus, Proteus, B. Cereus and sulfuric acid
bacteria). E.coli, S. Aureus, Proteus, B. Cereus and sulphite-reducing
clostridia) pathogenic microorganisms, including Salmonella, Listeria
monocytogenes; spoilage microorganisms.

Microbiological examination of the meat product allows us to assess its
condition and suitability for use in the production of meat products.

As a result of microbiological analysis the compliance of
microbiological indicators with the established indicators according to
the normative documentation was established. Coliforms, L.monocytogenes,
pathogenic microorganisms including salmonella were not detected.
Microbiological parameters of intestinal raw horse meat after ultrasonic
treatment are shown in Table 4. It can be seen that at 30 degrees all
microorganisms cease their activity.
\end{multicols}

\tcap{Table 3 - Organoleptic and physicochemical characteristics of horse intestinal raw materials}
\begin{longtblr}[
  label = none,
  entry = none,
  note{} = {\emph{Microbiological analysis after ultrasound treatment.} In order to
ensure product safety, a system of hazard analysis and critical control
points (HACCP) is in place at the enterprises.},
]{
  width = \linewidth,
  colspec = {Q[202]Q[371]Q[367]},
  cells = {c},
  cells = {font = \small},
  row{1} = {font=\bfseries\small},
  hlines,
  vlines,
}
Name of indicators             & Characteristics and standards                                    & Results                                                          \\
1                              & 2                                                                & 3                                                                \\
Appearance                     & Moderately moist, not contaminated with foreign matter           & Moderately moist, not contaminated with foreign matter           \\
Colour                         & Ccream, light beige, beige                                       & Beige                                                            \\
Odour                          & Characteristic of good quality products, with no offensive odour & Characteristic of good quality products, with no offensive odour \\
pH meat extract, not less      & No more than 6,2                                                 & 6,12                                                             \\
Mass fraction of fat           & -                                                                & 9,86                                                             \\
Ammonia nitrogen determination & negative reaction                                                & negative reaction                                                \\
Protein determination          & -                                                                & 9,035                                                            
\end{longtblr}

\tcap{Table 4 - Microbiological parameters of intestinal raw horse
meat after ultrasonic treatment}
\begin{longtblr}[
  label = none,
  entry = none,
]{
  width = \linewidth,
  colspec = {Q[542]Q[308]Q[90]},
  cells = {c},
  cells = {font = \small},
  row{1} = {font=\bfseries\small},
  hlines,
  vlines,
}
Name of indicators                                           & Characteristics and standards & Results                   \\
1                                                            & 2                             & 3                         \\
NMAFAnM, CFU/mL, not more                                    & 5х10\textsuperscript{6}       & 3,2х10\textsuperscript{4} \\
Presence of coliforms                                        & not allowed                   & absent                    \\
Presence of sulphite-reducing clostridia in 0.1 g of product & not allowed                   & absent                    \\
Pathogenic, including salmonellae in 25g                     & not allowed                   & absent                    
\end{longtblr}

\begin{multicols}{2}
As a result of microbiological analysis, compliance of microbiological
indicators with the established indicators according to the regulatory
documentation was established.

Microbiological analysis of intestinal raw materials showed that in the
process of heat treatment contamination of raw materials is
significantly reduced (from 260 000 to 600 microbes in 1 g of the
product), such microorganisms as salmonellae, proteins, sulfite
reducing, clostridia, coaguladepositive staphylococci and other
pathogenic microflora after heat treatment are not detected.

The table 5 presents the results of the safety assessment of horsemeat
intestinal by-products, including the levels of toxic elements (Pb, Cd,
Hg, As), organochlorine pesticides (HCH isomers, DDT and its
metabolites) and radionuclides (Cs-137, Sr-90). All indicators were
evaluated in accordance with the requirements of TR CU 021/2011 «On Food
Safety». Analytical measurements were carried out using standardized
methods: atomic absorption spectrometry for toxic elements (GOST
30178-96), gas-liquid chromatography for pesticide residues (GOST
23452-2015), and gamma-spectrometric analysis for radionuclides (GOST
32161-2013). The obtained results demonstrate that all contaminants were
absent or significantly below the permissible limits, confirming full
compliance of the products with food safety standards.
\end{multicols}

\tcap{Table 5 - Chemical, toxicological and radiological safety
indicators of horsemeat intestinal by-products in accordance with TR CU
021/2011 and relevant GOST standards}
\begin{longtblr}[
  label = none,
  entry = none,
]{
  width = \linewidth,
  colspec = {Q[215]Q[417]Q[108]Q[196]},
  cells = {c},
  cells = {font = \small},
  row{1} = {font=\bfseries\small},
  row{2} = {font=\bfseries\small},
  row{5} = {font=\bfseries\small},
  row{10} = {font=\bfseries\small},
  cell{2}{1} = {c=4}{0.936\linewidth},
  cell{5}{1} = {c=4}{0.936\linewidth},
  cell{10}{1} = {c=4}{0.936\linewidth},
  hlines,
  vlines,
}
Indicator                        & Regulatory limit (TR CU 021/2011 / GOST) & Result     & Compliance         \\
Pesticides, mg/kg, not more:     &                                          &            &                    \\
HCH (α, β, γ isomers)            & {≤ 0.1\\(GOST 23452-2015)}               & absent     & Meets requirements \\
DDT and metabolites              & {≤ 0.1\\(GOST 23452-2015)}               & absent     & Meets requirements \\
Toxic elements, mg/kg, not more: &                                          &            &                    \\
Lead (Pb)                        & {≤ 0.5\\(GOST 30178-96)}                 & absent     & Meets requirements \\
Arsenic (As)                     & {≤ 0.1\\(GOST 30178-96)}                 & absent     & Meets requirements \\
Cadmium (Cd)                     & {≤ 0.03\\(GOST 30178-96)}                & 0.03       & Meets requirements \\
Mercury (Hg)                     & {≤ 0.1\\(GOST 30178-96)}                 & absent     & Meets requirements \\
Radionuclides, Bq/kg, not more:  &                                          &            &                    \\
Cesium-137                       & {≤ 200\\(GOST 32161-2013)}               & 0.0 ± 12.2 & Below limit        \\
Strontium-90                     & no limit in TR CU                        & 0.0 ± 1.0  & Safe               
\end{longtblr}

\begin{multicols}{2}
According to the results of the analysis of food safety indicators the
meat product complies with the food safety norms, the content of
radionuclides - caesium-137 and strontium -- 90 does not exceed the
maximum permissible norm (200 bq/kg) and corresponds to 0,0\(\pm\)12,2;
0,0\(\pm\)1,0 bq/kg. All the results of the food safety analysis
indicate that the intestinal raw material after ultrasonic bath
treatment has no effect on the food safety indicators and is ready for
further sale.

Thus, the optimum ratio of processing of intestinal raw materials, which
have a positive effect on microbiological, physico-chemical and food
safety parameters, was identified. The research methods were carried out
according to the normative documents in the RSE "Centre of
Sanitary-Epidemiological Expertise" of the Medical Centre of the Medical
Centre of the Presidential Administration Medical Centre of the
Presidential Administration of the Republic of Kazakhstan" on the right
of economic management, Astana, Kazakhstan.

Comparative analysis of the treatment groups revealed distinct
differences in microbial dynamics and physicochemical stability. The
control samples demonstrated the fastest increase in TPC values,
exceeding 6.0 log CFU/g by day 9. Samples treated with citric acid
(0.5-1.0\%) showed a slower growth rate and maintained TPC below 5.0 log
CFU/g at the same storage point. Ultrasound treatment alone resulted in
an additional reduction of 0.3-0.5 log CFU/log compared with citric acid
treatment. The combined citric acid and ultrasound treatment
demonstrated the strongest effect, maintaining the lowest TPC values
(below 4.0 log CFU/log on day 9) and ensuring the most stable pH profile
throughout storage. These results confirm a synergistic antimicrobial
interaction between ultrasound and citric acid. The synergistic
reduction in microbial load is attributed to cavitation-induced
structural damage of bacterial cells and the acidification effect of
organic acids, which together enhance membrane permeability and
accelerate bacterial inactivation.

{\bfseries Conclusion.} The study showed that intestinal offal from horse
meat (``Qarta'') is very sensitive to microbiological contamination,
which limits their safe use as raw materials for the production of
delicatessen products. Microbiological assessment confirmed that
untreated intestines carry a high microbial load, while treatment with
ultrasound and citric acid significantly increases safety. Combined
treatment with 1.0\% citric acid (5 min) followed by ultrasound (35 kHz,
200 W, 5 min) was recognized as the most effective method to reduce the
total number of viable microorganisms from 2.6 × 10⁵ to 6.0 × 102 CFU/g
and maintain a low level of microbial growth during 13 days of storage.
in the fridge. After treatment, pathogenic microorganisms such as
Salmonella spp., Listeria monocytogenes and Escherichia coli were not
detected. In addition, the pH values remained stable within the
acceptable technological range (5.26--5.52), which confirms the
stability of the product. These results indicate that treatment with
citric acid using ultrasound can serve as a promising approach to ensure
microbiological safety and prolong the shelf life of horse intestines,
which will allow them to be used in the production of high-quality
delicatessen meat products. The combined use of ultrasound and mild
organic acid treatment can therefore be considered a promising and
technologically feasible approach for improving the microbiological
safety «Qarta» and similar intestinal by-products.
\end{multicols}

\begin{center}
{\bfseries Литература}
\end{center}

\begin{refs}
1. Rathnayake P.Y., Yu R., Yeo S.E., Choi Y.-S., Hwangbo S., Yong H.I.
Application of ultrasound to animal-based food to improve microbial
safety and processing efficiency//Food Science of Animal Resources.
-2025. -Vol.45(1). -P.199-222. DOI 10.5851/kosfa.2024.e128.

2. Bariya A.R., Rathod N.B., Patel A.S., Nayak J.K.B., Ranveer R.C.,
Hashem A., Abd\_Allah E.F., Ozogul F., Jambrak A.R., Rocha J.M. Recent
developments in ultrasound approach for preservation of animal origin
foods//Ultrasonics Sonochemistry. -2023. -Vol.101:106676. DOI
10.1016/j.ultsonch.2023.106676.

3. Chandrapala J., Oliver C., Kentish S., Ashokkumar M. Ultrasonics in
food processing //Ultrasonics Sonochemistry. -2012. -Vol.19(5). -P.
975--983. DOI 10.1016/j.ultsonch.2012.01.010.

4. Sagong H.-G., Lee S.-Y., Chang P.-S., Heu S., Ryu S., Choi Y.-J., Kang
D.-H. Combined effect of ultrasound and organic acids to reduce
Escherichia coli O157:H7, Salmonella Typhimurium, and Listeria
monocytogenes on organic fresh lettuce //International Journal of Food
Microbiology. -2011. -Vol.145(1). -P.287-292. DOI
10.1016/j.ijfoodmicro.2011.01.010.

5. Char C.D., Mitilinaki E., Guerrero S.N., Alzamora S.M. Use of
high-intensity ultrasound and UV-C light to inactivate some
microorganisms in fruit juices //Food and Bioprocess Technology.
-2010. -Vol.3(6). -P.797-803. DOI 10.1007/s11947-009-0307-7.

6. Gallo M., Ferrara L., Naviglio D. Application of ultrasound in food
science and technology: a perspective //Foods. -2018. -Vol.7
(10):164. DOI 10.3390/foods7100164.

7. Sorathiya K.B., Melo A., Hogg M.C., Pintado M. Organic acids in food
preservation: exploring synergies, molecular insights, and sustainable
applications //Sustainability. -2025. -Vol.17(8): 3434. DOI
10.3390/su17083434.

8. Zhang M., Haili N., Chen Q., Xia X., Kong B. Influence of
ultrasound-assisted immersion freezing on the freezing rate and
quality of porcine longissimus muscles// Meat Science. -2018. -Vol.
136. -P.1-8. DOI 10.1016/j.meatsci.2017.10.005.

9. Alarcon-Rojo A.D., Carrillo-Lopez L.M., Reyes-Villagrana R.,
Huerta-Jiménez M., Garcia-Galicia I.A. Ultrasound and meat quality: a
review //Ultrasonics Sonochemistry. -2019. -Vol.55. -P.369--382. DOI
10.1016/j.ultsonch.2018.09.016.

10. Faustman C., Sun Q., Mancini R., Suman S.P. Myoglobin and lipid
oxidation interactions: mechanistic bases and control //Meat Science.
-2010. -Vol.86(1). -P.86-94. DOI 10.1016/j.meatsci.2010.04.025.

11. Alarcon-Rojo A.D., Janacua H., Rodriguez J.C., Paniwnyk L., Mason T.J.
Power ultrasound in meat processing//Meat Science. -2015. -Vol.107.
-P.86-93. DOI 10.1016/j.meatsci.2015.04.015.

12. Kasaai M.R. Input power-mechanism relationship for ultrasonic
irradiation: food and polymer applications //Natural Science. -2013.
-Vol.5(8). -P.14-22. DOI 10.4236/ns.2013.58a2003.

13. Stanisławczyk R., Żurek J., Rudy M., Gil M. Influence of horse age on
carcass tissue composition and horsemeat quality: exploring
nutritional and health benefits for gourmets // Applied Sciences.
-2023. -Vol.13(20):11293. DOI 10.3390/app132011293.

14. Технический регламент Таможенного союза ТР ТС 021/2011. О безопасности
пищевой продукции. -Москва: Евразийская экономическая комиссия. -2011.
-46 с.

15. Продукты пищевые. Метод выявления бактерий рода Salmonella: ГОСТ
31659-2012 (ISO 6579:2002). -Москва: Госстандартинформ. - 2014. -25 с.

16. Продукты пищевые. Методы выявления и определения количества бактерий
группы кишечных палочек (колиформных бактерий): ГОСТ 31747-2012. -
Москва: Госстандартинформ. - 2013. -16 с.

17. Ben-David A., Davidson C.E. Estimation method for serial dilution
experiments// Journal of Microbiological Methods. -2014. -Vol.107.
-P.214--221. DOI 10.1016/j.mimet.2014.08.023.

18. Erkmen O. Gram staining technique // Laboratory Practices in
Microbiology. -2021. -P.99--105. DOI
10.1016/b978-0-323-91017-0.00028-7.

19. Kostanova A., Baytukenova S., Baytukenova S., Ryspaeva U., Yussupova
G., Shadyarova Z., Bekakhmetov A. Impact of combined ultrasound and
citric acid treatments on the quality and safety of natural horse
casings over storage durations // Polish Journal of Food and Nutrition
Sciences. -2025. -Vol 75(1). -P.37- 48. DOI 10.31883/pjfns/200346.

20. Siró I., Vén Cs., Balla Cs., Jónás G., Zeke I., Friedrich L.
Application of an ultrasonic assisted curing technique for improving
the diffusion of sodium chloride in porcine meat //Journal of Food
Engineering. -2009. -Vol.91(2). -P.35- 362. DOI
10.1016/j.jfoodeng.2008.09.015.
\end{refs}

\begin{center}
{\bfseries References}
\end{center}

\begin{refs}
1. Rathnayake P.Y., Yu R., Yeo S.E., Choi Y.-S., Hwangbo S., Yong H.I.
Application of ultrasound to animal-based food to improve microbial
safety and processing efficiency//Food Science of Animal Resources.
-2025. -Vol.45(1). -P.199-222. DOI 10.5851/kosfa.2024.e128.

2. Bariya A.R., Rathod N.B., Patel A.S., Nayak J.K.B., Ranveer R.C.,
Hashem A., Abd\_Allah E.F., Ozogul F., Jambrak A.R., Rocha J.M. Recent
developments in ultrasound approach for preservation of animal origin
foods//Ultrasonics Sonochemistry. -2023. -Vol.101:106676. DOI
10.1016/j.ultsonch.2023.106676.

3. Chandrapala J., Oliver C., Kentish S., Ashokkumar M. Ultrasonics in
food processing //Ultrasonics Sonochemistry. -2012. -Vol.19(5). -P.
975--983. DOI 10.1016/j.ultsonch.2012.01.010.

4. Sagong H.-G., Lee S.-Y., Chang P.-S., Heu S., Ryu S., Choi Y.-J.,
Kang D.-H. Combined effect of ultrasound and organic acids to reduce
Escherichia coli O157:H7, Salmonella Typhimurium, and Listeria
monocytogenes on organic fresh lettuce //International Journal of Food
Microbiology. -2011. -Vol.145(1). -P.287-292. DOI
10.1016/j.ijfoodmicro.2011.01.010.

5. Char C.D., Mitilinaki E., Guerrero S.N., Alzamora S.M. Use of
high-intensity ultrasound and UV-C light to inactivate some
microorganisms in fruit juices //Food and Bioprocess Technology. -2010.
-Vol.3(6). -P.797-803. DOI 10.1007/s11947-009-0307-7.

6. Gallo M., Ferrara L., Naviglio D. Application of ultrasound in food
science and technology: a perspective //Foods. -2018. -Vol.7 (10):164.
DOI 10.3390/foods7100164.

7. Sorathiya K.B., Melo A., Hogg M.C., Pintado M. Organic acids in food
preservation: exploring synergies, molecular insights, and sustainable
applications //Sustainability. -2025. -Vol.17(8): 3434. DOI
10.3390/su17083434.

8. Zhang M., Haili N., Chen Q., Xia X., Kong B. Influence of
ultrasound-assisted immersion freezing on the freezing rate and quality
of porcine longissimus muscles// Meat Science. -2018. -Vol.136. -P.
1-8. DOI 10.1016/j.meatsci.2017.10.005.

9. Alarcon-Rojo A.D., Carrillo-Lopez L.M., Reyes-Villagrana R.,
Huerta-Jiménez M., Garcia-Galicia I.A. Ultrasound and meat quality: a
review //Ultrasonics Sonochemistry. -2019. -Vol.55. -P.369--382. DOI
10.1016/j.ultsonch.2018.09.016.

10. Faustman C., Sun Q., Mancini R., Suman S.P. Myoglobin and lipid
oxidation interactions: mechanistic bases and control //Meat Science.
-2010. -Vol.86(1). -P.86-94. DOI 10.1016/j.meatsci.2010.04.025.

11. Alarcon-Rojo A.D., Janacua H., Rodriguez J.C., Paniwnyk L., Mason
T.J. Power ultrasound in meat processing//Meat Science. -2015. -Vol.
107. -P.86-93. DOI 10.1016/j.meatsci.2015.04.015.

12. Kasaai M.R. Input power-mechanism relationship for ultrasonic
irradiation: food and polymer applications //Natural Science. -2013.
-Vol.5(8). -P.14-22. DOI 10.4236/ns.2013.58a2003.

13. Stanisławczyk R., Żurek J., Rudy M., Gil M. Influence of horse age
on carcass tissue composition and horsemeat quality: exploring
nutritional and health benefits for gourmets // Applied Sciences. -2023.
-Vol.13(20):11293. DOI 10.3390/app132011293.

14. Tehnicheskij reglament Tamozhennogo sojuza TR TS 021/2011. O
bezopasnosti pishhevoj produkcii. -Moskva: Evrazijskaja jekonomicheskaja
komissija. -2011. -46 s. {[}in Russian{]}

15. Produkty pishhevye. Metod vyjavlenija bakterij roda Salmonella: GOST
31659-2012 (ISO 6579:2002). -Moskva: Gosstandartinform. - 2014. -25 s.
{[}in Russian{]}

16. Produkty pishhevye. Metody vyjavlenija i opredelenija kolichestva
bakterij gruppy kishechnyh palochek (koliformnyh bakterij): GOST
31747-2012. - Moskva: Gosstandartinform. - 2013. -16 s. {[}in Russian{]}

17. Ben-David A., Davidson C.E. Estimation method for serial dilution
experiments// Journal of Microbiological Methods. -2014. -Vol.107. -P.
214--221. DOI 10.1016/j.mimet.2014.08.023.

18. Erkmen O. Gram staining technique // Laboratory Practices in
Microbiology. -2021. -P.99--105. DOI
10.1016/b978-0-323-91017-0.00028-7.

19. Kostanova A., Baytukenova S., Baytukenova S., Ryspaeva U., Yussupova
G., Shadyarova Z., Bekakhmetov A. Impact of combined ultrasound and
citric acid treatments on the quality and safety of natural horse
casings over storage durations // Polish Journal of Food and Nutrition
Sciences. -2025. -Vol 75(1). -P.37- 48. DOI 10.31883/pjfns/200346.

20. Siró I., Vén Cs., Balla Cs., Jónás G., Zeke I., Friedrich L.
Application of an ultrasonic assisted curing technique for improving the
diffusion of sodium chloride in porcine meat //Journal of Food
Engineering. -2009. -Vol.91(2). -P.35- 362. DOI
10.1016/j.jfoodeng.2008.09.015.
\end{refs}

\begin{info}
\hspace{1em}\emph{{\bfseries Information about the authors}}

Baytukenova Sh. - candidate of technical sciences, associate professor,
S.Seifullin Kazakh Agrotechnical Research University, Astana,
Kazakhstan, e-mail:
baytukenova75@mail.ru;

Kostanova A. - PhD student, S.Seifullin Kazakh Agrotechnical Research
University, Astana, Kazakhstan, e-mail: anel\_kostanova@mail.ru;

Baytukenova S. - candidate of technical sciences, associate professor,
K.Kulazhanov Kazakh University of Technology and Business, Astana,
Kazakhstan, e-mail:
saule7272@mail.ru.

\hspace{1em}\emph{{\bfseries Сведения об авторах}}

Байтукенова Ш. - кандидат технических наук, ассоциированный профессор,
НАО «Казахский агротехнический исследовательский университет им.
С.Сейфуллина», Астана, Казахстан, e-mail: baytukenova75@mail.ru;

Костанова А.- докторант, НАО «Казахский агротехнический
исследовательский университет им. С.Сейфуллина», Астана, Казахстан,
e-mail: anel\_kostanova@mail.ru;

Байтукенова С. - кандидат технических наук, ассоциированный профессор,
АО «Казахский университет технологии и бизнеса им. К. Кулажанова»,
Астана, Казахстан, e-mail: saule7272@mail.ru.
\end{info}
